
\documentclass{report}

\input{../../preamble}
\input{../../macros}
\input{../../letterfonts}
\usepackage{parskip}
\title{HW04 Template}
\author{Ryan Lin}
\date{}

\begin{document}
\maketitle
\section*{Problem 1: Aces and Face Cards}

A standard deck consists of 52 cards of which 4 are aces, 4 are kings, and 12 (including the four kings) are "face cards" (Jacks, Queens, and Kings).

Cards are dealt at random without replacement from a standard deck till all the cards have been dealt. 


\noindent
\textbf{a)} Let $I_i$ be if card $i$ that is drawn is a face card, where $I_i \sim Bern(12/52)$ i.i.d.. Let $X$ be the number of face cards among the first 5 cards, thus $X = \sum_{i=1}^{5} I_i$

$$E[I_1] = 12/52$$

$$E[X] = \sum_{i=1}^5 E[I_i] = 5 \cdot 12/52$$


\noindent
\textbf{b.)} The number of face cards that *do not* appear among the first 13 cards


Let $I_i$ be if card $i$ that is drawn is a face card, where $I_i \sim Bern(12/52)$ i.i.d.. Let $X$ be the number of face cards that do not appear among the first 13 cards, thus $X = 12 - \sum_{i=1}^{13} I_i$

\begin{align*}
	E[X] &= E[12 - \sum_{i=1}^{13} I_i] \\
	&= 12 - \sum_{i=1}^{13}E[I_i] \\
	&= 12 - \sum_{i=1}^{13} \frac{12}{52} \\
	&= 12 - 13 \cdot \frac{12}{52} \\
	&= 9 \\
\end{align*}


\noindent
\textbf{c.)} The number of aces among the first 5 cards minus the number of kings among the last 5 cards

Let $ A_i$ be 1 if card $ i $ is an ace, otherwise 0. Let $ K_j $ be 1 if card $ j $ is a king, otherwise 0. $ A_j, K_j \sim \Bern(\frac{4}{52})  $ 

Let $ X $ be the number of aces among the first 5 cards minus the number of kings among the last 5 cards.
\begin{align*}
	E[X] &= E[\sum_{i=1}^{5} A_i - \sum_{j=48}^{52} K_j] \\
&= \sum_{i=1}^{5} E[A_i] - \sum_{j=48}^{52} K_j \\
&= 5 \cdot \frac{4}{52} - 5 \cdot \frac{4}{52} \\
&= 0 \\
\end{align*}


\noindent
\textbf{d.)} The number of cards before the first face card. 

If we take $ X_i $ to be the gap length between a face card.

Basically, we want to find the average number of cards between each gap. 

There are 13 possible gaps, with the total gap length of $ \frac{52-12} =40$, or there are 40 non face cards in the deck.  Since we know that the sample is a simple random sample, then we know that the length of the gaps come from a uniform distribution. So $ E[X_1] = E[X_2] = E[X_3] $ etc.

Thus $$ \sum_{i=1}^{13} E[X_i] = E[X_1] + E[X_2] + \dots + E[X_{13}] = 40 $$

So $ \boxed{E[X_1] = \frac{40}{13}} $ 




\noindent
\textbf{e.)} The number of cards strictly in between the first face card and the last face card

The number of cards between the first face card and the last face card is seeing how many cards are between the 2nd and 2nd to last gaps. Let $ A $ represent this number. 

In other words, we are finding the expected number of cards in 11 gaps, plus the 10 face cards between those gaps. From d, we knew the expected length of a gap was $E[X_i]= \frac{40}{13} $. 

Thus

$$  E[A] = \sum_{i=2}^{12} E[X_i] + E[10] = 11E[X_1] + 10 = 11 \cdot \frac{40}{13} + 10 $$

\noindent
\textbf{f.)} The number of aces before the first face card

We can also use the method of gaps here, where we can find the expected number of aces within one gap, where the gap is defined by the face card. 

Let $ A $ be the number of aces before the first face card. Let $ X $ be the number of cards before the first face card. 


We know that the number of cards before the first face card is a random variable. However, if it is not, the number of aces before the first face cards becomes a hypergeometric distribution problem.  $ A \mid X \sim \text{hypergeom}(40, 4, X) $

For example, 
$$P(A = a \mid X = x ) = \frac{\binom{4}{a}\binom{36}{x-4}}{\binom{40}{x}}$$

We know that the expected value of a hypergeometric variable is $ n \frac{G}{N} = X \frac{4}{40}$ 

\begin{align*}
E[A] &= E[E[A \mid X]] \\
&= E\left[ X \cdot \frac{4}{40} \right] \\
&= \frac{4}{40} \cdot E[X] \\
&= \frac{4}{40} \cdot \frac{40}{13} \\
&= \frac{4}{13} \\
\end{align*}

\hrulefill

\textbf{BAD PATHWAY} 

$ X $ is also a hyper geometric distribution, because we are trying to see the probability of not drawing any face cards. 

$$P(X=x) = \frac{\binom{40}{x}}{\binom{52}{x}}$$

\begin{align*}
	E[A] &= E[E[A \mid X]] \\
	&= \sum_{x=1}^{52} E[A \mid X=x] P(X=x)  \\
\end{align*}

We know that there are 12 face cards, therefore there is no way that there can be more than 40 cards before the first face cards. 


\begin{align*}
	&= \sum_{x=1}^{40} E[A \mid X=x] P(X=x) \\
	&= \sum_{x=1}^{40} x \cdot \frac{4}{52} \cdot P(X=x) \\
	&= \sum_{x=1}^{40} x \cdot \frac{4}{52} \cdot \frac{\binom{40}{x}}{\binom{52}{x}} \\
\end{align*}



\newpage

\section*{2.) Collecting Distributed Values}

This exercise is a workout in finding expectations by using all the tools at your disposal. If an answer doesn't appear to fit into a formula that has already been proven, it's a very good idea to try to write the variable as a sum of simpler variables.

**Notice that each part asks for the "expected number of" times something happens. But a different method is required for each part.**

\textbf{a)} A fair die is rolled $n$ times. Find the expected number of times the face with six spots appears.

Let $ I_i = \begin{cases}
	1, \quad \text{if side with 6 dots appears on roll $ i $ } \\
	0, \quad \text{otherwise}
\end{cases} \sim \Bern(\frac{1}{6}) $ 
$$E[X] = \sum_{i=1}^{n} E[I_i] = n\frac{1}{6}$$

\textbf{b}) A fair die is rolled $n$ times. Find the expected number of faces that *do not* appear, and say what happens to this expectation as $n$ increases.

Let $ I_i = \begin{cases}
	1, \text{if side with $ i $ dots does not appear} \\
	0, \text{otherwise}
\end{cases} \sim \Bern(p) $ i.i.d.

\[
p = \frac{5}{6}^{n}
\]
Thus, if we let $ X $ be the number of faces that do not appear: 
$$X =  \sum_{i=1}^{6} I_i$$

\begin{align*}
	E[X] &= E[\sum_{i=1}^{6}I_i] \\
	&= \sum_{i=1}^{6}E[I_i] \\
&=  \sum_{i=1}^{6} E[I_i] \\
&=  6E[I_1] \\ 
&= 6 \left( \frac{5}{6} \right) ^{n} \\
\end{align*}

As $n$ increases, the expected number of faces that do not appear decreases. This makes intuitive sense because a fair die draws from a uniform distribution. As the number of rolls increases, then it is less likely for one face not to show up. 

\textbf{c.)} Use your answer to Part **b** to find the expected number of distinct faces that *do* appear in $n$ rolls of a die.


\[
6-E[X] = 6- 6\left( \frac{5}{6} \right) ^{n}
\]

\textbf{d.)} Find the expected number of times you have to roll a die till you have seen all of the faces. This is a version of what is known as the *coupon collector's problem*. [For this one, as for many probability problems, it's a good idea to pretend you're doing the experiment. What happens on the first roll? What will you be waiting for next? What's the distribution of that waiting time? And so on.]


As you keep rolling the probability of getting a new face decreases. It's almost a geometric distribution, but the trials are not independent. There are 6 events that happen when the probability changes. First, when we get one face, then the probability to find a new face drops to $ \frac{5}{6} $. When we get the next, it drops to $ \frac{4}{6} $, etc. 

We can model this through 6 events then. 

Let $ X $ be the number of times you roll a die until you have seen all of the faces. 

Let $ X_i \sim \text{geometric}(p_i) $ be the number of times you roll until you have seen $ i $ faces. $ p_i = \frac{6-(i-1)}{6} $ 

\begin{align*}
	E[X] &= E\left[\sum_{i=1}^{6}X_i \right]\\
	&= \sum_{i=1}^{6}E[X_i] \\
	&= E[X_1] + E[X_2] + \dots + E[X_6]\\
	&= \frac{1}{1} + \frac{1}{\frac{5}{6}} + \dots + \frac{1}{\frac{1}{6}} \\
	&= 1 + \frac{6}{5} + \frac{6}{4} + \frac{6}{3} + \frac{6}{2} + \frac{6}{1} \\
	&= 14.7 \\
\end{align*}


\newpage
\section*{3. Unbiased Estimators} 

Remember that an estimator is a random variable that you must be able to calculate based on the observed data. That is, it must be a function of the observed variable or variables. Its formula cannot include the unknown parameter that you are trying to estimate. 

\textbf{a)} A population of known size $N$ contains an unknown number $G$ of good elements. Let $X$ be the number of good elements in a simple random sample of size $n$ drawn from this population. Use $X$ to construct an unbiased estimator of $G$, in the following steps:

1. Start with a known connection between $X$ and $G$: Write $E(X)$ in terms of $G$ and the known parameters.

2. Next, isolate $G$ in your equation above.

3. Use Step 2 as inspiration to identify a random variable whose expectation is $G$. [Remember: $X$ is a random variable. $E(X)$ isn't.]


$ X $ is a hypergeometric distribution with parameters $ X \sim \text{hypergeometric}(N, G, n) $ 

\begin{align*}
	E[X] &= n \frac{G}{N} \\
	G &= \frac{N\cdot E[X]}{n} \\
\end{align*}

Since we know that $ N $ is a constant and is known, as well a $ n $ is the sample size, these are both constants. Therefore, we can have random variable $ T = \frac{N}{n} X$ 

\[
	E[T] = E\left[\frac{N}{n} X \right] = \frac{N}{n}E[X] = \frac{N}{n} \cdot n \frac{G}{N} = G 
\]

\textbf{b)} Would your answer to Part **a** have been different if $X$ had been the number of good elements in a random sample drawn with replacement from the population? Why or why not?

No, even if the sample was drawn with replacement, we would have said that $ X \sim \Binom(n, \frac{G}{N}) $. This is because if we take getting a Good element as a success, we have $ n $ independent and identically distributed trials with success $ \frac{G}{N} $. We are then counting the number of successes in the sample.

We would then have:

\[
E[X] = np = n \frac{G}{N}
\]

Which is the exact same expression as before, which would result in the same unbiased estimator working. 



\textbf{c)} A flattened die lands 1 and 6 with chance $p/2$ each, and the other faces 2, 3, 4, and 5 with chance $(1-p)/4$ each. Here $p \in (0, 1)$ is an unknown number. Let $X_1, X_2, \ldots, X_n$ be the results of $n$ rolls of this die. First find $E(\vert X_1 - 3.5 \vert)$, and use the answer to construct an unbiased estimator of $p$ based on all of $X_1, X_2, \ldots, X_n$.

We can use the range space definition of expectation here. 

\begin{align*}
E[\vert X_1 - 3.5 \vert] &= (|1-3.5|) \cdot \frac{p}{2} + (|2-3.5|) (\frac{1-p}{4}) + (|3-3.5|) (\frac{1-p}{4}) + (|4-3.5|)(\frac{1-p}{4}) + (|5-3.5|)(\frac{1-p}{4}) + (|6-3.5|)(\frac{p}{2}) \\
&= (5)(\frac{p}{2}) + 4(\frac{1-p}{4})   \\
&= \frac{5}{2} p + (1-p) \\
&= 1 + \frac{3}{2}p \\
\end{align*}

Let $ \bar Y = \frac{1}{n} \sum_{i=1}^{n} |X_n-3.5| $ 

We proved in lecture that $ E[\bar Y] = E[|X_1 - 3.5|] $, because of symmetry and linearity of expectation. 

Therefore, 

\begin{align*}
	E[\bar Y] &= E[\vert X - 3.5 \vert ] \\ 
	&= 1 + \frac{3}{2} p \\
\end{align*}

Since an unbiased estimate would be $ E[T] = \theta $, we can to modify our current expression by a linear expression. 

\begin{align*}
	E[\frac{2}{3}(\bar Y-1)] &= \frac{2}{3} E[\bar Y] - \frac{2}{3} E[1] \\
	&= \frac{2}{3} (1+\frac{3}{2}p) - \frac{2}{3} \\
	&= \frac{2}{3} + p - \frac{2}{3} \\
	&=  p \\
\end{align*}

Thus, an unbiased estimator $ \boxed{T = \frac{2}{3} (\bar Y  - 1)}$

\newpage
\section*{4. Waiting for a Success
Run}\label{waiting-for-a-success-run}

\textbf{Required Reading:} Before you start this section, carefully read
\textbf{Answer 2} of
\href{https://data140.org/textbook/content/chapter-09/expected-waiting-times/\#waiting-till-hh}{Section
9.3.4}. That is the method you will apply in this exercise.

In a sequence of i.i.d. Bernoulli \((p)\) trials, let H represent
``heads'' or ``success'', that is, the event that occurs with
probability \(p\). Let \(W_{H,n}\) be the number of trials till you get
\(n\) heads in a row.

Remember that ``till'' means ``up to and including'' the \(n\) heads in
a row. Thus if the sequence of tosses starts out as TTTHTHHH then
\(W_{H,1} = 4\), \(W_{H,2} = 7\), and \(W_{H,3} = 8\).

\textbf{(a)} {[}These results are in the book but it will help to have
them here for reference. You don't have to prove them here.{]} Write
\(E(W_{H,1})\) and \(E(W_{H,2})\) as math expressions involving
\textbf{only} positive terms of the form \(1/p^k\) for positive integer
\(k\).

\begin{align*}
	E[W_{H,1}] &= \frac{1}{p} \\
	E[W_{H,2}] &= \frac{1}{p} + \frac{1}{p^{2}}
\end{align*}



\textbf{(b) Three heads in a row.} In Answer 2 of
\href{https://data140.org/textbook/content/chapter-09/expected-waiting-times/\#waiting-till-hh}{Section
9.3.4}, we condition on \(W_{H,1}\) to get \(E(W_{H,2})\). Now condition
on \(W_{H,2}\) to find \(E(W_{H,3})\). As above, your answer should
involve only positive terms of the form \(1/p^k\) for positive integer
\(k\). Find the numerical value of \(E(W_{H,3})\) in the case
\(p = 1/2\).


\begin{align*}
	E[W_{H,3}] &= E[W_{H,2}] + E[R] \\
	 &= E[W_{H,2}] + 1 + qE[W_{H,3}]\\
	 &= \frac{E[W_{H,2}] + 1}{1-q} \\
	 &= \frac{E[W_{H,2} +1]}{p} \\
	 &= \frac{\frac{1}{p} + \frac{1}{p^{2}}+1}{p} \\
	 &= \frac{1}{p} + \frac{1}{p^{2}} + \frac{1}{p^{3}} \\
\end{align*}


If $ p=\frac{1}{2} $ 

\[
E[H_{H,3}] = 2 + 4 + 8 = 14
\]




\textbf{(c) \(n\) heads in a row.} Guess an expression for
\(E(W_{H,n})\) for \(n \ge 1\). Prove it by induction, as follows:
Assume that your guess is true for \(n\), and show that then it must
also be true for \(n+1\).


Assume true for $ n $   
\[
E[W_{H,n}] = \sum_{i=1}^{n} \frac{1}{p^{i}}
\]

Show true for $ n+1 $ 

\begin{align*}
	E[W_{H,n+1}]&= E[H, n] + E[R] \\
	&= E[W_{H,n}] + 1 + qE[W_{H,n+1}] \\
	&= \frac{E[W_{H,n} + 1]}{1-q} \\
	&= \frac{E[W_{H,n} + 1]}{p} \\
	&= \frac{1}{p} + \frac{E[W_{H,n}]}{p} \\
	&= \frac{1}{p} + \frac{\sum_{i=1}^{n} \frac{1}{p^{i}}}{p} \\
	&= \frac{1}{p} + \sum_{i=1}^{n} \frac{1}{p^{i+1}} \\
	&= \frac{1}{p^{1}} + \sum_{i=2}^{n+1} \frac{1}{p^{i}} \\
	&=\boxed{ \sum_{i=1}^{n+1} \frac{1}{p^{i}} } \\
\end{align*}


\textbf{(d)} Define a function \texttt{ev\_W\_run} that takes \(p\) and
\(n\) as arguments and returns \(E(W_{H,n})\). Just replace the ellipsis
by an expression. Do not add any other lines of code.

\end{document}
